\section{Evaluation}\label{sec:eval}

In this section, we evaluate the effectiveness and efficiency of \tool. Our experiments are designed to answer the following three research questions:

\begin{itemize}
    \item \textbf{RQ1 (Effectiveness):} How effective is \tool in reproducing known consensus-breaking vulnerabilities compared to state-of-the-art baselines?
    \item \textbf{RQ2 (Ablation):} What is the contribution of each core design component (i.e., Specification-Guided Input Abstraction and KL-Divergence Reward) to the overall performance?
    \item \textbf{RQ3 (Zero-Day):} Is \tool capable of discovering previously unknown consensus-breaking vulnerabilities in production-grade clients?
\end{itemize}

\subsection{Experimental Setup}

\paragraph{Benchmarks and selection criteria}
% To evaluate the effectiveness of \tool, we compiled a benchmark suite of \tbd{XX} known consensus vulnerabilities and historical bugs from major Ethereum consensus clients (Prysm, Teku, Lighthouse, Nimbus) and core SSZ libraries. These benchmarks were selected based on three criteria: relevance (consensus-critical), reproducibility, and diversity (covering various bug types including logic errors, serialization edge cases, and panics).

To evaluate the effectiveness of \tool, we compiled a benchmark suite of \textbf{39} known SSZ vulnerabilities and bug instances drawn from production clients and SSZ libraries (see our public bug reports). The suite covers SSZ decoding, code generation, merkleization, HTR, proof logic, and schema-parameter validation across Go and Python implementations. All benchmarks are exercised through \tool's fuzzing pipeline; for codegen defects, a regeneration/build failure is treated as exposure and TTE is wall-clock time to the first failure.

\begin{table}
\centering
\caption{Benchmark sources (GitHub issues, 39 cases).}
\label{tab:bench-summary}
\small
\begin{tabular}{lccc}
\toprule
\textbf{Source} & \textbf{Lang} & \textbf{Count} & \textbf{Share} \\
\midrule
fastssz GitHub issues~\cite{fastssz} & Go & 29 & 74.36\% \\
py-ssz GitHub issues~\cite{pyssz} & Python & 10 & 25.64\% \\
\midrule
Total & - & 39 & 100.00\% \\
\bottomrule
\end{tabular}
\end{table}


\paragraph{Baselines}
We compare \tool against two state-of-the-art fuzzing approaches:
\begin{itemize}
    \item \textbf{BeaconFuzz:} The current industry-standard differential fuzzer for Ethereum 2.0~\cite{beaconfuzz}, which utilizes structure-aware mutation but lacks semantic guidance from the specification.
    \item \textbf{LibFuzzer (Coverage-Guided):} A general-purpose coverage-guided fuzzer~\cite{libfuzzer} without structural or specification guidance.
\end{itemize}
For full-suite reporting, we compare \tool directly against LibFuzzer and BeaconFuzz; both baselines are spec-blind, with BeaconFuzz adding structure-aware mutations.


\paragraph{Fuzzing policy initialization}
We employ two distinct initialization strategies for our experiments. For RQ1 and RQ2, to ensure a fair comparison with baselines, \tool starts with a \emph{randomly initialized policy} (cold start) for each benchmark run, learning the input distribution from scratch. However, for the RQ3 discovery study, we simulate a realistic deployment scenario by using a \emph{pre-trained policy} (warm start). This policy was pre-trained on the benchmark suite from RQ1, allowing the RL agent to transfer learned boundary-seeking strategies --- such as prioritizing offset edge cases or non-canonical padding --- to new, unseen targets.

\paragraph{Evaluation metrics}
We employ the following metrics to assess performance:
\begin{itemize}
    \item \textbf{Bug Reproduction Rate:} The number of benchmark bugs successfully triggered within the time limit.
    \item \textbf{Time to Exposure (TTE):} The time required to find the first failure-inducing input. TTE is marked as "--" if the bug was not found within the timeout.
    \item \textbf{Code Coverage:} Branch coverage achieved during the fuzzing session.
\end{itemize}

\paragraph{Environment}
All measurements were collected on a commodity workstation (Apple M3 Max, macOS 15.6.1) using Go 1.25.1. The default timeout for RQ1/RQ2 fuzzing sessions was 2h; the RQ3 zero-day campaign ran for 1 month.

\subsection{RQ1: Effectiveness on Benchmarks}

% We first evaluated whether \tool could reproduce the \tbd{XX} collected historical vulnerabilities. \autoref{tab:effectiveness} summarizes the results comparing \tool with BeaconFuzz and LibFuzzer.

We first evaluated whether \tool could reproduce the \textbf{39} collected benchmarks (29 fastssz regressions plus 10 py-ssz issues). \autoref{tab:effectiveness} reports the unified 39-case suite.


% \tool successfully reproduced \tbd{XX} out of \tbd{YY} benchmarks, significantly outperforming BeaconFuzz (which solved \tbd{ZZ}) and LibFuzzer (which solved \tbd{AA}). Notably, \tool demonstrated superior efficiency in detecting deep logic bugs. For example, in the case of a complex state transition bug in \textsc{Teku}, \tool found the violation in \tbd{XX} minutes, whereas BeaconFuzz required \tbd{YY} hours. The results indicate that \tool's strategy of leveraging specification-guided abstraction allows it to navigate the input space more effectively than coverage-guided baselines, particularly for vulnerabilities hidden behind strict validity checks.

\tool successfully reproduced all \textbf{39} benchmarks. LibFuzzer succeeds on 13/39 (33.33\%), while BeaconFuzz solves 20/39 (51.28\%) and \tool solves all 39. Across the suite, \tool exposes divergences quickly in a cold-start setting (median TTE 119.94\,ms), while the baselines are generally slower on boundary-heavy cases. These results indicate that \tool's strategy of leveraging specification-guided abstraction allows it to navigate the input space more effectively than coverage-guided baselines, particularly for vulnerabilities hidden behind strict validity checks.


\begin{table*}
\centering
\caption{Comparison of bug reproduction effectiveness on the 39-benchmark suite. TTE denotes Time to Exposure (lower is better). "--" indicates the bug was not found within the timeout.}
\label{tab:effectiveness}
\small
\rowcolors{4}{gray!10}{white}
\begin{tabular}{lp{2.1cm}p{7.4cm}rrr}
\toprule
\multirow{2}{*}{\textbf{Benchmark}} & \multirow{2}{*}{\textbf{SSZ Component}} & \multirow{2}{*}{\textbf{Bug Description}} & \multicolumn{3}{c}{\textbf{TTE (ms)}} \\
\cmidrule(lr){4-6}
 & & & \textbf{LibFuzzer} & \textbf{BeaconFuzz} & \textbf{\tool} \\
\midrule
FSSZ-181 & Codegen & codegen: fixed-size byte array size in external package & -- & -- & 1.64 \\
FSSZ-162 & Codegen & codegen: alias to non-\icode{uint64} integer types & -- & -- & 216.49 \\
FSSZ-152 & Codegen & codegen: unmarshal fixed-size custom types & -- & -- & 0.12 \\
FSSZ-127 & Codegen & codegen: incorrect output for nested \icode{[]byte} types & -- & -- & 0.19 \\
FSSZ-54 & Codegen & codegen: struct field not as pointer & -- & -- & 194.07 \\
FSSZ-76 & Codegen & codegen: alias/exclude-obj handling & -- & -- & 83.28 \\
FSSZ-153 & Codegen & codegen: cannot infer \icode{ssz-size} for fixed bytes & -- & -- & 223.69 \\
FSSZ-158 & Codegen & codegen: type alias on byte slice panics & -- & -- & 555.33 \\
FSSZ-166 & Codegen & codegen: external \icode{[]byte} import missing & -- & -- & 215.31 \\
FSSZ-136 & Codegen & codegen: unused imports emitted by generator & -- & -- & 206.61 \\
FSSZ-156 & Codegen & codegen: \icode{ssz-size} tag ignored with multiple tags & -- & -- & 81.76 \\
FSSZ-159 & Codegen & codegen: generic type parameters not resolved & -- & -- & 115.61 \\
FSSZ-164 & Codegen & codegen: external const array length unsupported & -- & -- & 80.00 \\
FSSZ-188 & Codegen & codegen: multiple field names in struct declaration & -- & -- & 86.59 \\
FSSZ-86 & Codegen & codegen: embedded container structs not expanded & -- & -- & 88.75 \\
FSSZ-100 & Codegen & codegen: missing time.Time support & -- & -- & 90.16 \\
FSSZ-149 & Codegen & codegen: list of \icode{uint32} unsupported & -- & -- & 93.45 \\
FSSZ-151 & Codegen & codegen: list of non-pointer containers unsupported & -- & -- & 203.28 \\
FSSZ-1 & Encoding & encoding: SizeSSZ panic on nil pointer fields & 1.46 & 0.11 & 0.07 \\
FSSZ-52 & Codegen & codegen: panic on \icode{ast.Interface} & -- & -- & 1.67 \\
FSSZ-173 & Proofs & proofs: skip intermediate hashes in multi-proof & -- & 232.62 & 146.74 \\
FSSZ-147 & Merkleization & merkleize: incorrect chunk count & -- & 40.76 & 26.70 \\
FSSZ-119 & Proofs & proofs: memory out of bounds & 14.56 & 0.79 & 127.78 \\
FSSZ-111 & Merkleization & merkleize: depth calculation for basic types & -- & 12.94 & 85.00 \\
FSSZ-110 & Merkleization & merkleize: generated hash digest not deterministic & 1371.18 & 1166.58 & 1334.89 \\
FSSZ-98 & HTR & HTR: allocate nil data structures & 1.58 & 0.10 & 0.11 \\
FSSZ-96 & HTR & HTR: allocate uint slice when len is 0 & -- & 30.89 & 93.36 \\
FSSZ-9 & Decoding & decode: bitlist parse panics on empty byte array & 9.88 & 0.18 & 0.24 \\
FSSZ-23 & Decoding & decode: incorrect offset handling in spec tests & 14.66 & 0.73 & 96.64 \\
PSSZ-109 & Decoding & decode: bitlist length underflow in \icode{Bitlist[2048]} & 130.73 & 163.22 & 207.35 \\
PSSZ-35 & Merkleization & merkleize: \icode{ByteList[31]} without padding & -- & 115.95 & 294.41 \\
PSSZ-74 & Decoding & decode: fixed-size list length check for \icode{BeaconBlockHeader[4]} & 242.79 & 109.33 & 314.29 \\
PSSZ-82 & HTR & HTR: bitlist length limit not enforced & 114.48 & 109.11 & 174.69 \\
PSSZ-83 & HTR & HTR: list length limit not enforced (\icode{List[uint64,128]}) & -- & 125.92 & 310.96 \\
PSSZ-113 & HTR & HTR: list missing \icode{mix\_in\_length} & -- & 127.16 & 355.61 \\ 
PSSZ-114 & HTR & HTR: bitlist missing \icode{mix\_in\_length} & 174.56 & 117.78 & 211.61 \\ 
PSSZ-111 & Schema validation & schema validation: accept \icode{Vector} length 0 & 128.78 & 115.43 & 119.94 \\
PSSZ-112 & Schema validation & schema validation: accept \icode{Bitvector} length 0 & 130.34 & 122.90 & 191.41 \\
PSSZ-116 & Schema validation & schema validation: reject \icode{List} max length 0 & 122.29 & 153.16 & 131.70 \\ 
\midrule
\rowcolor{gray!10}
 &  & \hfill \textbf{Avg. TTE (ms)} & \textbf{190.00} & \textbf{140.00} & \textbf{170.00} \\ 
 \rowcolor{gray!10}
 &  & \hfill \textbf{\#Solved} & \textbf{13} & \textbf{20} & \textbf{39} \\ 
 \rowcolor{gray!10}
 &  & \hfill \textbf{Success Rate} & \textbf{33.33\%} & \textbf{51.28\%} & \textbf{100.00\%} \\ 
\bottomrule
\end{tabular}
\end{table*}


% \begin{table}
% \centering
% \caption{Summary of effectiveness results on the 39-benchmark suite. Avg.\ TTE excludes timeouts.}
% \label{tab:results-summary}
% \small
% \begin{tabular}{lccc}
% \toprule
% \textbf{Fuzzer} & \textbf{Avg.\ TTE} & \textbf{Solved} & \textbf{Timeout} \\
% \midrule
% LibFuzzer & 0.19s & 13 (33.33\%) & 26 (66.67\%) \\ 
% BeaconFuzz & 0.14s & 20 (51.28\%) & 19 (48.72\%) \\ 
% \tool & 0.17s & 39 (100.00\%) & 0 (0.00\%) \\ 
% \bottomrule
% \end{tabular}
% \end{table}

\begin{mdframed}[
    backgroundcolor=gray!10,
    linecolor=white,
    skipabove=0.5em
    ]\noindent
    \textbf{Result for RQ1:} \tool reproduces all 39 benchmarks and maintains fast exposure times (median TTE 119.94\,ms), outperforming both baselines on success rate.
\end{mdframed}

\subsection{RQ2: Ablation Study}

To verify the contribution of our core design choices --- specifically the specification-guided input abstraction (SGIA) and the KL-divergence based reward mechanism --- we conducted an ablation study. We compared the full \tool implementation against two variants:
\begin{itemize}
    \item \textbf{\tool-NoSpec:} This variant disables the SGIA module. It initializes the fuzzing campaign with random valid seeds rather than boundary-seeking seeds derived from the specification.
    \item \textbf{\tool-NoRL:} This variant retains the SGIA but replaces the KL-divergence driven RL policy with a uniform random selection strategy for choosing abstract buckets.
\end{itemize}

\autoref{tab:ablation} presents the performance of these variants on the 39-case benchmark suite. Because Avg.\ TTE excludes timeouts, the primary signal is the number of bugs solved. \tool-NoSpec solved only 26 bugs (13 timeouts at the 2h limit); the similar Avg.\ TTE among solved cases (0.16s) does not compensate for the 13 unsolved benchmarks. \tool-NoRL matched the full system in success rate (39 bugs) but was slower overall (avg.\ TTE 0.28s vs.\ 0.17s; +58.82\%), showing the KL-driven policy improves efficiency even when SGIA remains intact.



\begin{table}
\centering
\caption{Comparison between ablative versions of \tool.}
\label{tab:ablation}
\small
\begin{tabular}{lccc}
\toprule
\textbf{Configuration} & \textbf{Bugs Found} & \textbf{Avg. TTE (ms)} & \textbf{$\Delta$ Avg. TTE vs Full} \\
\midrule
\tool-NoSpec & 26 & 160 & -5.26\% \\ 
\tool-NoRL & 39 & 280 & +58.82\% \\ 
\textbf{\tool (Full)} & \textbf{39} & \textbf{170} & -- \\ 
\bottomrule
\end{tabular}
\end{table}

\begin{figure}[]
    \footnotesize
    \centering
    \begin{tikzpicture}
        \begin{axis}[
            legend style={at={(0.09,0.7)}, anchor=west, nodes={scale=0.8, transform shape}, fill=white!90!gray, draw=none, rounded corners=3pt},
            legend cell align={left},
            width=\linewidth,
            height=2in, 
            xlabel=Number of Benchmarks Solved,
            ylabel=Cumulative Time (s),
            grid=major,
            minor grid style={gray!25}, 
            grid style={dashed, gray!20}, 
            tick label style={font=\small}, 
            mark size=1.0pt,
            line width=0.8pt,
            xmin=0,
            xmax=39,
            ymin=0,
            ymax=11, 
            ytick={0, 2, 4, 6, 8, 10}, 
            xtick={0, 10, 20, 30, 39},
            axis line style={thick},
            grid=major,
            grid style={dashed, gray!15},
            minor grid style={gray!10},
            font=\footnotesize,
            axis line style={black, line width=0.3pt},
            tick label style={font=\scriptsize},
            xlabel style={font=\sffamily, yshift=5pt},
            ylabel style={font=\sffamily, yshift=-10pt},
        ]
        \axispath\draw;
        % full version
        \addplot+[color=bluegray, thin, mark=*] plot coordinates {
(0, 0.00)
(1, 0.00)
(2, 0.00)
(3, 0.00)
(4, 0.00)
(5, 0.00)
(6, 0.00)
(7, 0.00)
(8, 0.03)
(9, 0.11)
(10, 0.19)
(11, 0.28)
(12, 0.36)
(13, 0.45)
(14, 0.54)
(15, 0.63)
(16, 0.72)
(17, 0.81)
(18, 0.91)
(19, 1.03)
(20, 1.15)
(21, 1.27)
(22, 1.40)
(23, 1.55)
(24, 1.73)
(25, 1.92)
(26, 2.11)
(27, 2.31)
(28, 2.52)
(29, 2.73)
(30, 2.94)
(31, 3.16)
(32, 3.37)
(33, 3.60)
(34, 3.89)
(35, 4.20)
(36, 4.52)
(37, 4.87)
(38, 5.43)
(39, 6.76)
        };

        % baseline
        \addplot+[color=amethyst, thin, mark=triangle*] plot coordinates {
(0, 0.00)
(1, 0.00)
(2, 0.00)
(3, 0.00)
(4, 0.00)
(5, 0.00)
(6, 0.00)
(7, 0.00)
(8, 0.08)
(9, 0.16)
(10, 0.25)
(11, 0.33)
(12, 0.42)
(13, 0.51)
(14, 0.61)
(15, 0.72)
(16, 0.85)
(17, 0.98)
(18, 1.12)
(19, 1.31)
(20, 1.52)
(21, 1.72)
(22, 1.94)
(23, 2.16)
(24, 2.38)
(25, 2.94)
(26, 4.27)
        };

        % baseline (FF)
        \addplot+[color=darkpastelgreen, thin, mark=diamond*] plot coordinates {
(0, 0.00)
(1, 0.00)
(2, 0.00)
(3, 0.00)
(4, 0.00)
(5, 0.00)
(6, 0.00)
(7, 0.00)
(8, 0.08)
(9, 0.17)
(10, 0.25)
(11, 0.34)
(12, 0.42)
(13, 0.51)
(14, 0.61)
(15, 0.70)
(16, 0.79)
(17, 0.89)
(18, 1.01)
(19, 1.13)
(20, 1.27)
(21, 1.41)
(22, 1.59)
(23, 1.79)
(24, 1.99)
(25, 2.19)
(26, 2.40)
(27, 2.61)
(28, 2.83)
(29, 3.06)
(30, 3.35)
(31, 3.66)
(32, 4.01)
(33, 4.38)
(34, 4.75)
(35, 5.25)
(36, 5.81)
(37, 7.14)
(38, 8.51)
(39, 10.74)
                 };
    
        \legend{\tool\\\tool-NoSpec\\\tool-NoRL\\}
    
        \end{axis}
    \end{tikzpicture}
    \caption{A comparison between \tool and its baselines without guidance from specification (\tool-NoSpec) and RL (\tool-NoRL) on the 39-case suite, where x-axis denotes the total number of benchmarks solved, and y-axis denotes the cumulative time spent in seconds.}
    \label{fig:ablation}
\end{figure}

\begin{mdframed}[
    backgroundcolor=gray!10,
    linecolor=white,
    skipabove=0.5em
    ]\noindent
    \textbf{Result for RQ2:} The primary gain is in success rate: SGIA raises solved benchmarks from 26 to 39 (\tool-NoSpec times out on 13/39). With SGIA fixed, the KL-driven policy improves efficiency, reducing average TTE from 0.28s to 0.17s.
\end{mdframed}

\subsection{RQ3: Zero-Day Vulnerability Discovery}

Finally, we evaluated \tool's capability to discover unknown vulnerabilities in the latest stable versions of production Ethereum consensus clients (Prysm~\cite{prysm}, Teku~\cite{teku}, Lighthouse~\cite{lighthouse}, Nimbus~\cite{nimbus}, Lodestar~\cite{lodestar}) and core SSZ libraries. Unlike the controlled experiments in RQ1, here \tool utilized the pre-trained policy derived from the benchmark suite, enabling it to apply learned boundary-seeking strategies to new targets immediately.

Over a period of \textbf{1} month, \tool discovered \textbf{12} previously unknown vulnerabilities. %of which \textbf{5} was confirmed as a consensus-breaking issue. 
These findings have been responsibly disclosed to the respective client teams. We highlight three representative case studies below.

\yanju{add 0day illustrative figures.}

\autoref{tab:0day-instances-report} lists the previously unknown vulnerability instances uncovered in this study. Each row corresponds to a distinct vulnerable implementation.

\begin{table*}[t]
\centering
\caption{Previously unknown vulnerability instances discovered by \tool in production SSZ implementations.}
\label{tab:0day-instances-report}
\small
\begin{tabular}{cll}
\toprule
\textbf{\#} & \textbf{Vulnerability} & \textbf{Implementation / Issue} \\
\midrule
1 & Container deserialize short/empty input & Lodestar~\cite{lodestar} \\
2.1 & Bitvector dirty padding & Teku~\cite{teku} \href{https://github.com/Consensys/teku/pull/10173}{Consensys/teku\#10173} \\
2.2 & Bitvector dirty padding & Prysm (fastssz)~\cite{prysm,fastssz} \\
2.3 & Bitvector dirty padding & ethereum/py-ssz~\cite{pyssz} \\
2.4 & Bitvector dirty padding & dynamic-ssz~\cite{dynamic_ssz} \href{https://github.com/pk910/dynamic-ssz/issues/54}{pk910/dynamic-ssz\#54} \\
3.1 & Boolean canonicalization \{0x00,0x01\} & fastssz / Prysm~\cite{fastssz,prysm} \href{https://github.com/ferranbt/fastssz/issues/222}{ferranbt/fastssz\#222} \\
3.2 & Boolean canonicalization \{0x00,0x01\} & dynamic-ssz~\cite{dynamic_ssz} \href{https://github.com/pk910/dynamic-ssz/issues/53}{pk910/dynamic-ssz\#53} \\
4.1 & Union trailing data accepted & Lighthouse~\cite{lighthouse} \href{https://github.com/sigp/ethereum_ssz/issues/64}{sigp/ethereum\_ssz\#64} \\
4.2 & Union trailing data accepted & Lodestar~\cite{lodestar} \\
4.3 & Union trailing data accepted & ethereum/remerkleable~\cite{remerkleable} \\
5 & Variable-length container gap & ethereum/remerkleable~\cite{remerkleable} \\
6 & Null-bitlist trap & dynamic-ssz~\cite{dynamic_ssz} \href{https://github.com/pk910/dynamic-ssz/issues/56}{pk910/dynamic-ssz\#56} \\
\bottomrule
\end{tabular}
\end{table*}


\subsubsection{Case Study 1: The Bitvector Dirty Padding Bug}

This vulnerability, found in Prysm~\cite{prysm} and \icode{py-ssz}~\cite{pyssz}, involves the improper validation of padding bits in \icode{Bitvector[N]} types. The SSZ specification mandates that unused high-order bits in the final byte must be zero. \tool synthesized a test case where the \icode{JustificationBits} (a \icode{Bitvector[4]}) had the dirty value \icode{0x1F} (binary \icode{0001 1111}) instead of the canonical \icode{0x0F}. \tool's RL policy successfully prioritized the \irule{NON\_CANONICAL\_SAME\_VALUE} abstract bucket, causing the concretizer to flip the padding bits. Vulnerable clients accepted this encoding, while strict clients like \textsc{Lighthouse}~\cite{lighthouse} rejected it, leading to a consensus split.

\subsubsection{Case Study 2: The Null-Bitlist Trap}

\tool discovered a critical flaw in the \icode{dynamic-ssz}~\cite{dynamic_ssz} library regarding \icode{Bitlist} deserialization. The specification requires a ``sentinel bit'' (1) to be appended to the bitstream to delimit its length; thus, an empty bitlist must be serialized as \icode{0x01}. \tool generated a boundary case containing a raw \icode{0x00} byte for a bitlist field. While a compliant decoder should reject this as missing the sentinel, the vulnerable implementation treated it as a valid empty list. This bug was uncovered because \tool explicitly modeled the "missing sentinel" boundary in its abstraction, guiding the fuzzer to generate inputs that coverage-guided tools typically filter out as invalid.

\subsubsection{Case Study 3: Variable-Length Container Gap}

In the \icode{remerkleable}~\cite{remerkleable} library, \tool identified a vulnerability where the decoder failed to enforce strict contiguity of offsets in variable-length containers. The spec requires that ${\sf Offset}_n = {\sf Offset}_{n-1} + {\sf Length}_{n-1}$. \tool synthesized a ``gap attack'' by generating a container where offsets were shifted to skip over 4 bytes of garbage data (``ghost bytes'') inserted between fields. The decoder simply followed the offsets, ignoring the gap, and successfully deserialized the object. However, this resulted in a different Merkle root compared to the canonical serialization. This finding demonstrates \tool's ability to exploit structural invariants that are often overlooked by developers focusing solely on field-level validation.

\begin{mdframed}[
    backgroundcolor=gray!10,
    linecolor=white,
    skipabove=0.5em
    ]\noindent
    \textbf{Result for RQ3:} \tool is highly effective at discovering zero-day vulnerabilities in production clients. Leveraging a pre-trained policy, it identified a critical consensus-breaking bug, demonstrating its practical value for securing live blockchain networks.
\end{mdframed}
